\section{Lecture 5}

\subsection*{Stationarity of a Nonlinear AR(1) Process}
Consider the stochastic recurrence equation

\[X_t = F(X_{t-1},\epsilon_t),\]

where \((\epsilon_t)_{t\in\mathbb Z}\) is an iid sequence.

We can think of the construction of a stationary solution as a stochastic analogue of the classical fixed-point theorem. To clarify, we are not solving

\[x^*=F(x^*,\epsilon)\]

for a point \(x^*\) in the state space. Rather, we seek a random variable \(X\), independent of a fresh innovation \(\epsilon\), such that

\[X\overset d=F(X,\epsilon).\]

For each innovation \(\epsilon\), define the stochastic Lipschitz constant

\[L_\epsilon
=
\sup_{x\neq x'}
\frac{|F(x,\epsilon)-F(x',\epsilon)|}{|x-x'|}.\]

If \(F(\cdot,\epsilon)\) is differentiable, then

\[L_\epsilon=\sup_x\left|\frac{\partial}{\partial x}F(x,\epsilon)\right|.\]

\subsection*{Theorem}
Suppose

\[\rho:=EL_\epsilon<1\]

and, for some \(x_0\in\mathbb R\),

\[E|F(x_0,\epsilon)|<\infty.\]

Then:

\begin{enumerate}
\item
  The stochastic recurrence equation admits a unique stationary causal solution
  \[X_t^*=g(\epsilon_t,\epsilon_{t-1},\ldots).\]
\item
  Its marginal distribution
  \[\pi=\mathcal L(X_t^*)\]
  is the unique stationary distribution among distributions with finite first moment.
\item
  For any deterministic initial value \(x\),
  \[X_n(x)\Rightarrow\pi.\]
\item
  More strongly, the process satisfies geometric moment contraction:
  \[E|X_n(x)-X_n(x')|
  \le
  \rho^n|x-x'|.\]
\end{enumerate}

\begin{center}\rule{0.5\linewidth}{0.5pt}\end{center}

\subsection*{Proof}
First observe that the integrability assumption at one point implies integrability at every point. For any \(x\),

\[\begin{aligned}
E|F(x,\epsilon)|
&\le
E|F(x_0,\epsilon)|
+
E|F(x,\epsilon)-F(x_0,\epsilon)|\\
&\le
E|F(x_0,\epsilon)|
+
EL_\epsilon|x-x_0|\\
&=
E|F(x_0,\epsilon)|+\rho|x-x_0|
<\infty.
\end{aligned}\]

Define the forward iteration starting from \(x\) by

\[X_n(x)
=
F_{\epsilon_n}\circ\cdots\circ F_{\epsilon_1}(x),
\qquad
F_\epsilon(x):=F(x,\epsilon).\]

Instead of studying the forward iteration directly, define the backward iteration

\[Z_n(x)
=
F_{\epsilon_0}\circ F_{\epsilon_{-1}}
\circ\cdots\circ F_{\epsilon_{-n+1}}(x).\]

Since the innovations are iid,

\[Z_n(x)\overset d=X_n(x).\]

We will show that \(Z_n(x)\) converges almost surely and in \(L^1\) to a random variable \(Z_\infty\) that does not depend on \(x\).

\subsection*{Step 1: Convergence of the backward iteration}
Consider successive backward iterations. We have

\[\begin{aligned}
|Z_n(x)-Z_{n-1}(x)|
&\le
L_{\epsilon_0}L_{\epsilon_{-1}}\cdots
L_{\epsilon_{-n+2}}
\,|F_{\epsilon_{-n+1}}(x)-x|.
\end{aligned}\]

Taking expectations and using independence,

\[E|Z_n(x)-Z_{n-1}(x)|
\le
\rho^{n-1}E|F(x,\epsilon)-x|.\]

Therefore

\[\sum_{n=1}^\infty
E|Z_n(x)-Z_{n-1}(x)|
<\infty.\]

By Tonelli's theorem,

\[E\left[
\sum_{n=1}^\infty
|Z_n(x)-Z_{n-1}(x)|
\right]
<\infty.\]

Hence

\[\sum_{n=1}^\infty
|Z_n(x)-Z_{n-1}(x)|
<\infty
\qquad\text{a.s.}\]

Thus \(Z_n(x)\) is almost surely Cauchy and therefore

\[Z_n(x)\to Z_\infty(x)
\qquad\text{a.s.}\]

Moreover, for \(m>n\),

\[E|Z_m(x)-Z_n(x)|
\le
\sum_{j=n+1}^mE|Z_j(x)-Z_{j-1}(x)|
\to0.\]

Hence \(Z_n(x)\) is Cauchy in \(L^1\), so

\[Z_n(x)\to Z_\infty(x)
\qquad\text{in }L^1.\]

\subsection*{Step 2: The limit does not depend on the initial condition}
For two initial conditions \(x,x'\), coupling the backward iterations using the same innovations gives

\[|Z_n(x)-Z_n(x')|
\le
\left(
\prod_{j=0}^{n-1}L_{\epsilon_{-j}}
\right)|x-x'|.\]

Taking expectations,

\[E|Z_n(x)-Z_n(x')|
\le
\rho^n|x-x'|.\]

Since both backward iterations converge in \(L^1\),

\[E|Z_\infty(x)-Z_\infty(x')|
=
0.\]

Therefore

\[Z_\infty(x)=Z_\infty(x')
\qquad\text{a.s.}\]

Thus the limiting random variable is independent of the initial condition. We denote it simply by \(Z_\infty\).

\subsection*{Step 3: Causal representation}
Each \(Z_n(x)\) is measurable with respect to

\[\sigma(\epsilon_0,\epsilon_{-1},\ldots,\epsilon_{-n+1}).\]

Since \(Z_n(x)\to Z_\infty\) almost surely,

\[Z_\infty
\in
\sigma(\epsilon_0,\epsilon_{-1},\epsilon_{-2},\ldots).\]

Hence there exists a measurable function \(g\) such that

\[Z_\infty
=
g(\epsilon_0,\epsilon_{-1},\epsilon_{-2},\ldots).\]

Define

\[X_t^*
=
g(\epsilon_t,\epsilon_{t-1},\ldots).\]

Because the innovation sequence is iid, \((X_t^*)\) is stationary. Furthermore, passing to the limit in the backward recursion gives

\[X_t^*=F(X_{t-1}^*,\epsilon_t)
\qquad\text{a.s.}\]

Thus \(X_t^*\) is a stationary causal solution. Let

\[\pi=\mathcal L(X_t^*).\]

\begin{center}\rule{0.5\linewidth}{0.5pt}\end{center}

\subsection*{Geometric Moment Contraction}
Consider two forward trajectories driven by the same future innovations:

\[X_n(x)
=
F_{\epsilon_n}\circ\cdots\circ F_{\epsilon_1}(x)\]

and

\[X_n(x')
=
F_{\epsilon_n}\circ\cdots\circ F_{\epsilon_1}(x').\]

Pathwise,

\[|X_n(x)-X_n(x')|
\le
\left(\prod_{j=1}^nL_{\epsilon_j}\right)|x-x'|.\]

Therefore

\[\boxed{
E|X_n(x)-X_n(x')|
\le
\rho^n|x-x'|.
}\]

This property is called \textbf{geometric moment contraction (GMC)}.

Now take \(X_0^*\sim\pi\), constructed from the past innovations

\[\epsilon_0,\epsilon_{-1},\ldots.\]

Then \(X_0^*\) is independent of the future innovations \(\epsilon_1,\ldots,\epsilon_n\). Couple the stationary and nonstationary chains using the same future innovations. Then

\[\begin{aligned}
E|X_n(x)-X_n(X_0^*)|
&\le
E\left[
\left(\prod_{j=1}^nL_{\epsilon_j}\right)
|x-X_0^*|
\right]\\
&=
\rho^nE|x-X_0^*|.
\end{aligned}\]

Since

\[X_n(X_0^*)\sim\pi,\]

we obtain

\[W_1\big(\mathcal L(X_n(x)),\pi\big)
\le
\rho^nE|x-X_0^*|.\]

Thus the distribution of the nonstationary chain converges exponentially fast to the stationary distribution in \(W_1\), and hence

\[X_n(x)\Rightarrow\pi.\]

\begin{center}\rule{0.5\linewidth}{0.5pt}\end{center}

\subsection*{\texorpdfstring{\(L^p\) Extension}{L\^{}p Extension}}
More generally, suppose for some \(p>0\),

\[\rho_p:=EL_\epsilon^p<1\]

and

\[E|F(x_0,\epsilon)|^p<\infty\]

for some \(x_0\).

Then

\[E|X_n(x)-X_n(x')|^p
\le
\rho_p^n|x-x'|^p.\]

Similarly,

\[E|Z_n(x)-Z_{n-1}(x)|^p
\le
C\rho_p^{n-1}.\]

For \(p\ge1\), the previous argument can be repeated using the \(L^p\) norm and Minkowski's inequality.

\textbf{Homework} Extend the proof to the case where \(0 < p < 1\).

\begin{center}\rule{0.5\linewidth}{0.5pt}\end{center}

\subsection*{Negative Lyapunov Exponent Condition}
The moment condition

\[EL_\epsilon^p<1\]

is stronger than what is needed merely for pathwise contraction.

A more general condition is

\[E\log L_\epsilon<0,\]

together with an appropriate integrability condition such as

\[E\log^+|F(x_0,\epsilon)-x_0|<\infty.\]

Indeed,

\[\frac1n
\log\left(
\prod_{j=1}^nL_{\epsilon_j}
\right)
=
\frac1n\sum_{j=1}^n\log L_{\epsilon_j}
\longrightarrow
E\log L_\epsilon<0
\qquad\text{a.s.}\]

by the strong law of large numbers. Hence

\[\prod_{j=1}^nL_{\epsilon_j}\]

decays exponentially fast almost surely.

The quantity

\[E\log L_\epsilon\]

can therefore be interpreted as a Lyapunov exponent for the stochastic dynamical system. A negative Lyapunov exponent says that trajectories driven by the same innovation sequence forget their initial conditions exponentially fast almost surely.

Notice that

\[EL_\epsilon<1
\quad\Longrightarrow\quad
E\log L_\epsilon
\le
\log EL_\epsilon<0\]

by Jensen's inequality. Thus the negative Lyapunov exponent condition is more general than first-moment contraction.

However,

\[E\log L_\epsilon<0\]

does not in general imply \(EL_\epsilon^p<1\) for a prescribed \(p\), so pathwise exponential contraction should be distinguished from \(p\)th geometric moment contraction.

\begin{center}\rule{0.5\linewidth}{0.5pt}\end{center}

\subsection*{Remark: Backward Iteration and Coupling from the Past}
The construction above is often called the \textbf{backward iteration} method for stochastic recurrence equations or iterated random functions.

It is closely related conceptually to coupling from the past, but it should be distinguished from the exact CFTP algorithm of Propp and Wilson.

Here we establish asymptotic coalescence:

\[|Z_n(x)-Z_n(x')|\to0
\qquad\text{as }n\to\infty.\]

In coupling from the past, one instead seeks a finite random time \(T\) such that trajectories from all possible initial states have exactly coalesced by time \(0\). This permits exact simulation from the stationary distribution after finite computation.

Ordinary Lipschitz contraction generally gives asymptotic coalescence, but not finite-time exact coalescence.
